\begingroup
\renewcommand{\arraystretch}{1.45}
\fancyhead[L]{\small TRANSCENDENTAL FUNCTIONS}
\Needspace{10\baselineskip}
\AGMArticle{11}
Now also easily the subsequent coefficients of the series

\par\noindent\hfill\rule{0.25\linewidth}{0.4pt}\hfill\mbox{}\par
\clearpage\selectlanguage{english}\fancyhead[L]{\small ARITHMETIC–GEOMETRIC MEAN.}
\pdfbookmark[0]{CONTINUATION OF THE INVESTIGATIONS ON THE ARITHMETIC–GEOMETRIC MEAN}{agm-fortsetzung}
\SourceLeaf{384}
\begin{center}\bfseries
CONTINUATION OF THE INVESTIGATIONS

ON THE ARITHMETIC–GEOMETRIC MEAN.

\rule{0.08\linewidth}{0.3pt}
\end{center}
[\textsc{Gauss}'s further investigations on the arithmetic–geometric mean were written down so incompletely, only in their final formulas, and with such changing notation, that almost all sense of the whole would be lost if the printed text reproduced only those formulas without accompanying explanation. In the following articles I have highlighted the passages found among the posthumous papers and connected them by a train of thought which may have prompted \textsc{Gauss} to set out the results in question from a recognizably common point of view.]

\Needspace{43\baselineskip}
\AGMArticle{12}
[In several places, besides the two original progressions of quantities $(a,b)$ considered in the preceding articles, which refer to the same arithmetic–geometric mean, \textsc{Gauss} also uses a third progression of quantities $(c)$. It is connected with the former by the equation $aa=bb+cc$, which holds for every common index of the three terms $(a,b,c)$. According to the conventions in Article 2, the equations for calculating the quantities derived from $a,b,c$ in the backward and forward directions can be put in the form:
\SourceLeaf{385}
\Needspace{10\baselineskip}
\[
\begin{aligned}&.\ .\quad{}^{\prime\prime}a={}^\prime a+{}^\prime c,\quad{}^\prime a=a+c,\quad a,\\&\qquad2a'=a+b,\quad2a''=a'+b',\quad.\ .\\[9pt]&.\ .\quad{}^{\prime\prime}b={}^\prime a-{}^\prime c,\quad{}^\prime b=a-c,\quad b,\\&\qquad b'b'=ab; \quad b''b''=a'b',\quad.\ .\\[9pt]&.\ .\quad{}^{\prime\prime}c{}^{\prime\prime}c=4{}^\prime a{}^\prime c,\quad{}^\prime c{}^\prime c=4ac,\quad c,\\&\qquad2c'=a-b,\quad2c''=a'-b',\quad.\ .\end{aligned}
\]
\Needspace{18\baselineskip}
From this it follows directly that, while
\Needspace{10\baselineskip}
\[
\begin{aligned}&\mathrm{M}(a,b)=\mathrm{M}(a^{\mathrm{n}},b^{\mathrm{n}})=\mathrm{M}({}^{\mathrm{n}}a,{}^{\mathrm{n}}b)\\&=\lim a^{\mathrm{m}}=\lim b^{\mathrm{m}}\quad\hbox{for}\quad m=\infty\end{aligned}
\]
holds,
\Needspace{10\baselineskip}
\[
\begin{aligned}&\mathrm{M}(a,c)=2^n\mathrm{M}(a^{\mathrm{n}},c^{\mathrm{n}})=2^{-n}\mathrm{M}({}^{\mathrm{n}}a,{}^{\mathrm{n}}c)\\&=\lim\frac{{}^{\mathrm{m}}a}{2^m}=\lim\frac{{}^{\mathrm{m}}c}{2^m}\quad\hbox{for}\quad m=\infty\end{aligned}
\]
holds, and that the remark appended to the first example in Article 3, that ${}^{\mathrm{n}}a,{}^{\mathrm{n}+1}a,{}^{\mathrm{n}+2}a$ increase approximately like the terms of a geometric progression with the ratio $2$, is thus shown to hold generally; further, that for the case calculated in Example 4 of Article 3, $a=b\surd2=c\surd2$, the following hold for every $n$: ${}^{\mathrm{n}}a=2^n.a^{\mathrm{n}}$, ${}^{\mathrm{n}}b=2^n.c^{\mathrm{n}}$, ${}^{\mathrm{n}}c=2^n.b^{\mathrm{n}}$.

\Needspace{12\baselineskip}
\par The way in which these quantities approach their limits follows from the following series expansions
\Needspace{10\baselineskip}
\[
\begin{aligned}\mathrm{M}(a,b)&=a^{\mathrm{n}}-c^{\mathrm{n}+1}-c^{\mathrm{n}+2}\\&\quad-.\;.\;.-c^{\mathrm{n}+\mathrm{m}}-.\;.\;.\end{aligned}
\]
\Needspace{10\baselineskip}
\[
\begin{aligned}\mathrm{M}(a,b)&=b^{\mathrm{n}}+c^{\mathrm{n}+1}-c^{\mathrm{n}+2}\\&\quad-.\;.\;.-c^{\mathrm{n}+\mathrm{m}}-.\;.\;.\end{aligned}
\]
\Needspace{10\baselineskip}
\[
\begin{aligned}\mathrm{M}(a,c)&=2^{-n}.{}^{\mathrm{n}}a-2^{-n-1}.{}^{\mathrm{n}+1}b-2^{-n-2}.{}^{\mathrm{n}+2}b\\&\quad-.\;.\;.-2^{-n-m}.{}^{\mathrm{n}+\mathrm{m}}b-.\;.\;.\end{aligned}
\]
\Needspace{10\baselineskip}
\[
\begin{aligned}\mathrm{M}(a,c)&=2^{-n}.{}^{\mathrm{n}}c+2^{-n-1}.{}^{\mathrm{n}+1}b-2^{-n-2}.{}^{\mathrm{n}+2}b\\&\quad-.\;.\;.-2^{-n-m}.{}^{\mathrm{n}+\mathrm{m}}b-.\;.\;.\end{aligned}
\]
\Needspace{10\baselineskip}
\[
\begin{aligned}\mathrm{M}(a,b)^2&=a^{\mathrm{n}}.a^{\mathrm{n}}-\frac12c^{\mathrm{n}}.c^{\mathrm{n}}-\frac32c^{\mathrm{n}+1}.c^{\mathrm{n}+1}\\&\quad-.\;.\;.-\frac32c^{\mathrm{n}+\mathrm{m}}.c^{\mathrm{n}+\mathrm{m}}-.\;.\;.\end{aligned}
\]
\Needspace{10\baselineskip}
\[
\begin{aligned}\mathrm{M}(a,b)^2&=b^{\mathrm{n}}.b^{\mathrm{n}}+\frac12c^{\mathrm{n}}.c^{\mathrm{n}}-\frac32c^{\mathrm{n}+1}.c^{\mathrm{n}+1}\\&\quad-.\;.\;.-\frac32c^{\mathrm{n}+\mathrm{m}}.c^{\mathrm{n}+\mathrm{m}}-.\;.\;.\end{aligned}
\]
\Needspace{10\baselineskip}
\[
\begin{aligned}\mathrm{M}(a,c)^2&=2^{-2n}.{}^{\mathrm{n}}c.{}^{\mathrm{n}}c+\frac12 2^{-2n}.{}^{\mathrm{n}}b.{}^{\mathrm{n}}b-\frac32 2^{-2n-2}.{}^{\mathrm{n}+1}b.{}^{\mathrm{n}+1}b\\&\quad-.\;.\;.\qquad-\frac32 2^{-2n-2m}.{}^{\mathrm{n}+\mathrm{m}}b.{}^{\mathrm{n}+\mathrm{m}}b-.\;.\;.\end{aligned}
\]
\Needspace{10\baselineskip}
\[
\begin{aligned}\mathrm{M}(a,c)^2&=2^{-2n}.{}^{\mathrm{n}}a.{}^{\mathrm{n}}a-\frac12 2^{-2n}.{}^{\mathrm{n}}b.{}^{\mathrm{n}}b-\frac32 2^{-2n-2}.{}^{\mathrm{n}+1}b.{}^{\mathrm{n}+1}b\\&\quad-.\;.\;.\qquad-\frac32 2^{-2n-2m}.{}^{\mathrm{n}+\mathrm{m}}b.{}^{\mathrm{n}+\mathrm{m}}b-.\;.\;.\end{aligned}
\]
\Needspace{10\baselineskip}
\[
\begin{aligned}\surd\mathrm{M}(a,b)&=\surd a^{\mathrm{n}}-\surd c^{\mathrm{n}+2}-\surd c^{\mathrm{n}+4}\\&\quad-.\;.\;.-\surd c^{\mathrm{n}+2\mathrm{m}}-.\;.\;.\end{aligned}
\]
\Needspace{10\baselineskip}
\[
\begin{aligned}\surd\mathrm{M}(a,b)&=\surd b^{\mathrm{n}}+\surd c^{\mathrm{n}+2}-\surd c^{\mathrm{n}+4}\\&\quad-.\;.\;.-\surd c^{\mathrm{n}+2\mathrm{m}}-.\;.\;.\end{aligned}
\]
\Needspace{12\baselineskip}
\SourceLeaf{386}
\Needspace{10\baselineskip}
\[
\begin{aligned}\surd\mathrm{M}(a,c)&=\surd2^{-n}.{}^{\mathrm{n}}a-\surd2^{-n-2}.{}^{\mathrm{n}+2}b-\surd2^{-n-4}.{}^{\mathrm{n}+4}b\\&\quad-.\;.\;.-\surd2^{-n-2m}.{}^{\mathrm{n}+2\mathrm{m}}b-.\;.\;.\end{aligned}
\]
\Needspace{10\baselineskip}
\[
\begin{aligned}\surd\mathrm{M}(a,c)&=\surd2^{-n}.{}^{\mathrm{n}}c+\surd2^{-n-2}.{}^{\mathrm{n}+2}b-\surd2^{-n-4}.{}^{\mathrm{n}+4}b\\&\quad-.\;.\;.-\surd2^{-n-2m}.{}^{\mathrm{n}+2\mathrm{m}}b-.\;.\;.\end{aligned}
\]
\Needspace{10\baselineskip}
\[
\begin{aligned}\displaystyle\frac\pi2.\frac{\mathrm{M}(a,b)}{\mathrm{M}(a,c)}&\displaystyle=\frac1{2^n}\log4\frac{a^{\mathrm{n}}}{c^{\mathrm{n}}}-\frac1{2^n}\log\frac{a^{\mathrm{n}}}{a^{\mathrm{n}+1}}\\&\quad-.\;.\;.-\frac1{2^{n+m}}\log\frac{a^{\mathrm{n}+\mathrm{m}}}{a^{\mathrm{n}+\mathrm{m}+1}}-.\;.\end{aligned}
\]
\Needspace{10\baselineskip}
\[
\begin{aligned}\displaystyle\frac\pi2.\frac{\mathrm{M}(a,b)}{\mathrm{M}(a,c)}&\displaystyle=\frac1{2^n}\log4\frac{b^{\mathrm{n}+1}}{c^{\mathrm{n}}}+\frac3{2^{n+1}}\log\frac{b^{\mathrm{n}+2}}{b^{\mathrm{n}+1}}\\&\quad+.\;.\;.+\frac3{2^{n+m}}\log\frac{b^{\mathrm{n}+\mathrm{m}+1}}{b^{\mathrm{n}+\mathrm{m}}}+.\;.\end{aligned}
\]
\Needspace{10\baselineskip}
\[
\begin{aligned}\displaystyle\frac\pi2.\frac{\mathrm{M}(a,c)}{\mathrm{M}(a,b)}&\displaystyle=\frac1{2^n}\log2\frac{{}^{\mathrm{n}}a}{{}^{\mathrm{n}}b}-\frac1{2^n}\log2\frac{{}^{\mathrm{n}}a}{{}^{\mathrm{n}+1}a}\\&\quad-.\;.\;.-\frac1{2^{n+m}}\log2\frac{{}^{\mathrm{n}+\mathrm{m}}a}{{}^{\mathrm{n}+\mathrm{m}+1}a}-.\;.\end{aligned}
\]
\Needspace{10\baselineskip}
\[
\begin{aligned}\displaystyle\frac\pi2.\frac{\mathrm{M}(a,c)}{\mathrm{M}(a,b)}&\displaystyle=\frac1{2^n}\log2\frac{{}^{\mathrm{n}+1}c}{{}^{\mathrm{n}}b}+\frac3{2^{n+1}}\log\frac12\frac{{}^{\mathrm{n}+2}c}{{}^{\mathrm{n}+1}c}\\&\quad+.\;.\;.+\frac3{2^{n+m}}\log\frac12\frac{{}^{\mathrm{n}+\mathrm{m}+1}c}{{}^{\mathrm{n}+\mathrm{m}}c}+.\;.\end{aligned}
\]
\Needspace{27\baselineskip}
\par In the last four equations, $\frac\pi2$ has been put in place of the limits, for continually increasing $m$, of respectively
\Needspace{10\baselineskip}
\[
\begin{gathered}\frac{\mathrm{M}(a^{\mathrm{m}},c^{\mathrm{m}})}{\mathrm{M}(a^{\mathrm{m}},b^{\mathrm{m}})}.\log4\frac{a^{\mathrm{m}}}{c^{\mathrm{m}}};\\[9pt]
\frac{\mathrm{M}(a^{\mathrm{m}},c^{\mathrm{m}})}{\mathrm{M}(a^{\mathrm{m}},b^{\mathrm{m}})}.\log4\frac{b^{\mathrm{m}+1}}{c^{\mathrm{m}}},\\[9pt]
\frac{\mathrm{M}({}^{\mathrm{m}}a,{}^{\mathrm{m}}b)}{\mathrm{M}({}^{\mathrm{m}}a,{}^{\mathrm{m}}c)}.\log4\frac{{}^{\mathrm{m}}a}{{}^{\mathrm{m}}b},\\[9pt]
\frac{\mathrm{M}({}^{\mathrm{m}}a,{}^{\mathrm{m}}b)}{\mathrm{M}({}^{\mathrm{m}}a,{}^{\mathrm{m}}c)}.\log2\frac{{}^{\mathrm{m}+1}c}{{}^{\mathrm{m}}b}\end{gathered}
\]
\Needspace{25\baselineskip}
that is, in place of the limit of $\mathrm{M}(1,\varepsilon)\log\frac4\varepsilon$ for positive values of $\varepsilon$ decreasing to zero. For it follows first from those equations, in which, as the relations
\Needspace{10\baselineskip}
\[
\begin{gathered}\frac{a'}a=1-\frac{c'}a,\\[7pt]
\left(\frac{b'}b\right)^4=1+\frac{cc}{bb},\\[7pt]
\frac{{}^\prime a}a=2-\frac{{}^\prime b}a,\\[7pt]
\left(\frac12\frac{{}^\prime c}c\right)^4=1+\frac{bb}{cc}\end{gathered}
\]
readily show, the values whose logarithms are to be taken, as soon as they are all real, decrease steadily to unity from the second term of the progression onward, that the quantity sought is fully determined and, for example, if $a=b\surd2=c\surd2$ is set, is a quantity lying between the bounds $\frac94\log2$ and $\frac{10}4\log2$. The above equations also hold for complex values of the variables if $n=0$ is set and those values of $\log{}^{\mathrm{m}}a$, $\log{}^{\mathrm{m}}c$, $\log a$, $\log b$, $\log c$, $\log a^{\mathrm{m}}$, $\log b^{\mathrm{m}}$ are used which follow by continuous change from the real quantities. Now take as one system $(a,b,c)$ the system consisting of real quantities and subject to the condition $a=b\surd2=c\surd2$, and as another system $(\alpha,\beta,\gamma)$ that for which $\alpha=c$, $\beta=b\surd{-1}$, $\gamma=a=b\surd2=c\surd2$ and the values of $\beta^{\mathrm{m}}$ and ${}^{\mathrm{m}}\gamma$ to be determined by extracting roots have positive real parts. A third variable system $(A,B,C)$ can then be set up which passes continuously from the one $(a,b,c)$ to the other $(\alpha,\beta,\gamma)$, and does so
\SourceLeaf{387}
in such a way that during this passage none of the variables ${}^{\mathrm{m}}A$, ${}^{\mathrm{m}}C$, $A$, $B$, $C$, $A^{\mathrm{m}}$, $B^{\mathrm{m}}$ touches the value zero or acquires a negative real part. The last of those equations, with $n$ set equal to zero, therefore holds both for the system $(a,b,c)$ and for $(\alpha,\beta,\gamma)$, and since
\Needspace{10\baselineskip}
\[
\begin{gathered}\mathrm{M}(\alpha,\beta)=\tfrac12(1+\surd{-1})\mathrm{M}(a,b),\\[7pt]
\mathrm{M}(\alpha,\gamma)=\mathrm{M}(a,c)=\mathrm{M}(a,b)\end{gathered}
\]
holds, $\frac{\pi}{2}\surd{-1}=\log\surd{-1}$ follows; thus $\pi$ has its usual meaning as the ratio of the length of the circumference of a circle to its diameter.

\Needspace{16\baselineskip}
\par With the aid of this consideration, the correctness of the following remark recorded by \textsc{Gauss} can also be established. I give it the notation used here and follow it with an example calculated by \textsc{Gauss}:]
\par \textit{The arithmetic–geometric means take a different form if the negative value is chosen for one of $b', b'', b'''\ .\ .$: yet all the results are included in the following form:}
\Needspace{10\baselineskip}
\[
\frac{1}{\mathfrak{M}(a,b)}=\frac{1}{\mathrm{M}(a,b)}+\frac{4ik}{\mathrm{M}(a,c)}
\]
[where $k$ denotes a real integer.]

\Needspace{38\baselineskip}
\par \textit{Example of an imaginary value of the A. G. mean}
\Needspace{10\baselineskip}
\[
\begin{array}[t]{r@{\,=\,}r l r}
a&3.0000000&\log\ .\ .&0.4771213\\
b&1.0000000&&0.0000000\\[6pt]
a'&2.0000000&&0.3010300\\
b'&1.7320508&&0.2385606\\[6pt]
a''&1.8660254&&0.2709175\\
b''&1.8612098&&0.2697953\\[6pt]
a'''&1.8636176&&0.2703568\\
b'''&1.8636159&&0.2703564\\[6pt]
a''''&1.8636167&&0.2703566\\\hline
\end{array}
\]
\Needspace{10\baselineskip}
\[
\begin{array}[t]{r@{\,=\,}l l r}
a&3.0000000&\log\ .\ .&0.4771213\\
c&2.8284270&&0.4515450\\[6pt]
\frac12\mathbin{.}{}'a&2.9142\mathrm{i}35&&0.4645214\\[6pt]
\frac12\mathbin{.}{}'c&2.9129510&&0.4643332\\[6pt]
\frac14\mathbin{.}{}''a&2.9135822&&0.4644273\\[4pt]\hline
\end{array}
\]
\Needspace{43\baselineskip}
\SourceLeaf{388}
\begingroup\small\setlength{\arraycolsep}{2pt}
\Needspace{10\baselineskip}
\[
\begin{array}{r@{\,=\,}r@{}l r r r r}
a&3.0000000&&\log\ .\ .\ 0.4771213&0&&\\
b&1.0000000&&0.0000000&360^{0}&&\\[6pt]
a'&2.0000000&&0.3010300&0&&\\
b'&-1.7320508&&0.2385606&180^{0}&&\\[6pt]
a''&0.1339746&&9.1270225&0&&\\
b''&&{}+1.8612098\,i&0.2697953&90^{0}&&\\[6pt]
a'''&0.0669873&{}+0.9306049\,i&9.9698876&85^{0}&52'&58''10\\
b'''&0.3530969&{}+0.3530969\,i&9.6984089&45&0&0\\[6pt]
a''''&0.2100421&{}+0.6418509\,i&9.8295254&71^{0}&52&46.58\\
b''''&0.2836930&{}+0.6208239\,i&9.8341482&65&26&29.05\\[6pt]
a^{\mathrm{V}}&0.2468676&{}+0.6313374\,i&9.8311572&68&38&36.05\\
b^{\mathrm{V}}&0.2470649&{}+0.6324002\,i&9.8318368&68&39&37.82\\[6pt]
a^{\mathrm{VI}}&0.24699625&{}+0.6318688\,i&9.8314971&68&39&6.95\\
b^{\mathrm{VI}}&0.24699625&{}+0.6318685\,i&9.8314970&68&39&6.93\\[6pt]
a^{\mathrm{VII}}&0.24699625&{}+0.63186865\,i&9.83149705&68&39&6.94\\\hline
\end{array}
\]
\endgroup
\Needspace{10\baselineskip}
\[
\begin{aligned}&\frac{1}{\mathfrak{M}(a,b)}=+0.5365910-1.3728774\,i\\&=\frac{1}{\mathrm{M}(a,b)}+\frac{4i}{\mathrm{M}(a,c)}\end{aligned}
\]
\clearpage
\endgroup
